\section{引言}

\subsection{研究背景}

人工智能应用的进入门槛持续降低，公开网络中出现了大量与 AI、Agent、RAG 和自动化有关的项目、讨论及资讯。然而，记录中同时混有终端用户的求解诉求、开发者内部任务、产品介绍、新闻转述和营销表达。若将这些文本直接按出现频率或平台排序汇总，容易把技术热度误写成真实需求，把同一事件的转载误写成多源验证，也容易忽略正文不完整和热度指标缺失带来的不确定性。

赛题要求系统围绕给定需求方向完成多源采集清洗、机会判别与分类、主题聚合、价值评估和决策支持五项任务。其难点不在于生成一个单独的分类标签，而在于把原始来源、正文证据、机会条件、主题成员和价值分数串成可追溯链路。本文将这条链路作为研究对象，研究问题包括：如何在来源异质、内容完整度不同的情况下保留可解释的数据证据；如何用结构化机会条件区分需求与宣传；如何在确认边界不足时把模型输出与代理复核清晰分层；以及如何防止主题聚合和价值排序超出数据实际支持范围。

\subsection{研究内容与贡献}

本文的工作集中在以下三个方面。第一，设计了批次隔离的多源数据处理流程，保存原始字段、清洗记录、正文快照、去重关系和质量标记。第二，建立了带证据回溯的机会判别与分层复核流程。模型判断由可观察需求、AI 核心性、技术可实现性和排除项共同决定，复核结果区分原模型、可见标签辅助复核、盲复核和 AI 裁定，避免把模型标签写成真人真值。第三，构建语义主题和价值排序的探索性计算，并以审计指标明确主题稀疏、来源偏差和缺少人工评价时的结论边界。

本文不把 B01 公开数据当作所有平台的完整样本，也不把代理复核当作人工标注。论文的结论主要回答系统是否形成了可复现、可解释的机会发现流程，以及当前样本能支持哪些探索性观察。

研究问题据此分为三层：多源数据在怎样的保留条件下能够形成可审计的样本；结构化机会条件及多轮模型对照能揭示哪些判断边界；语义聚类和加权评分在何种情况下只适合作为人工筛选工具。本文将来源覆盖、内容质量和复核状态一起报告，使主题排序能够连回其样本范围。
